\documentclass[11pt]{article}
\usepackage[margin=1in]{geometry}
\usepackage{amsmath}
\usepackage{amssymb}
\usepackage{xcolor}
\usepackage{tikz}

\title{Spring 2025 MC Problem 3: Perfect Complements\\Marshallian vs. Hicksian Demand}
\author{ECO 310}
\date{}

\begin{document}

\maketitle

\section{The Problem}

Assume utility $u(x_1, x_2) = \min\{4x_1, x_2\}$. The Hicksian demand bundle is given by:

\begin{enumerate}
    \item[(a)] $(4u^*, u^*)$
    \item[(b)] $\left(\frac{u^*}{4}, u^*\right)$ \quad \textcolor{red}{$\checkmark$ CORRECT}
    \item[(c)] $\left(u^*, \frac{u^*}{4}\right)$
    \item[(d)] $(u^*, 4u^*)$
\end{enumerate}

The problem asks for \textbf{Hicksian demand}, but let's compare it with Marshallian demand to understand the difference.

\section{Perfect Complements Intuition}

With $u(x_1, x_2) = \min\{4x_1, x_2\}$, goods are consumed in a \textbf{fixed ratio} to avoid waste.

\subsection{The Kink Condition}

Utility is determined by the minimum of $4x_1$ and $x_2$. To maximize utility without wasting money:
\[
\boxed{4x_1 = x_2}
\]

This is like left shoes and right shoes: you need them in a specific ratio. Here, for every 1 unit of good 1, you need 4 units of good 2.

\subsection{Graphical Representation}

\begin{center}
\begin{tikzpicture}[scale=1.2]
% Axes
\draw[->] (0,0) -- (5,0) node[right] {$x_1$};
\draw[->] (0,0) -- (0,5) node[above] {$x_2$};

% Indifference curves (L-shaped)
\draw[thick, blue] (0.5, 5) -- (0.5, 2) -- (5, 2);
\node[blue, right] at (5, 2) {$u = 2$};

\draw[thick, blue] (1, 5) -- (1, 4) -- (5, 4);
\node[blue, right] at (5, 4) {$u = 4$};

% Optimal consumption line
\draw[thick, red, dashed] (0, 0) -- (4, 4);
\node[red, above right] at (2.5, 2.5) {$x_2 = 4x_1$};

% Kink points
\filldraw[blue] (0.5, 2) circle (2pt);
\filldraw[blue] (1, 4) circle (2pt);

% Labels
\node[below] at (0.5, 0) {$0.5$};
\node[below] at (1, 0) {$1$};
\node[left] at (0, 2) {$2$};
\node[left] at (0, 4) {$4$};

\end{tikzpicture}
\end{center}

The consumer always chooses bundles along the line $x_2 = 4x_1$ (the red dashed line where the kinks occur).

\section{Marshallian Demand (Income Fixed)}

\subsection{The Problem}

\textbf{Maximize} utility subject to a budget constraint:
\[
\max_{x_1, x_2} \, \min\{4x_1, x_2\} \quad \text{subject to} \quad p_1 x_1 + p_2 x_2 = m
\]

\textbf{What's fixed:} Income $m$ and prices $(p_1, p_2)$

\textbf{Question:} ``Given my income $m$, what should I buy?''

\subsection{Solution}

\textbf{Step 1:} Set the kink condition $4x_1 = x_2$ (optimal allocation)

\textbf{Step 2:} Substitute into the budget constraint:
\begin{align*}
p_1 x_1 + p_2 x_2 &= m \\
p_1 x_1 + p_2(4x_1) &= m \\
x_1(p_1 + 4p_2) &= m
\end{align*}

\textbf{Step 3:} Solve for $x_1$ and $x_2$:
\[
\boxed{x_1^*(p_1, p_2, m) = \frac{m}{p_1 + 4p_2}}
\]
\[
\boxed{x_2^*(p_1, p_2, m) = 4x_1^* = \frac{4m}{p_1 + 4p_2}}
\]

\subsection{Properties}

\begin{itemize}
    \item Depends on \textbf{income $m$} and \textbf{prices}
    \item Observable in real markets
    \item Also called ``uncompensated demand''
\end{itemize}

\section{Hicksian Demand (Utility Fixed)}

\subsection{The Problem}

\textbf{Minimize} expenditure subject to achieving a target utility:
\[
\min_{x_1, x_2} \, p_1 x_1 + p_2 x_2 \quad \text{subject to} \quad \min\{4x_1, x_2\} = u^*
\]

\textbf{What's fixed:} Target utility $u^*$ and prices $(p_1, p_2)$

\textbf{Question:} ``What's the cheapest way to achieve utility $u^*$?''

\subsection{Solution}

\textbf{Step 1:} Set the kink condition $4x_1 = x_2$ (optimal allocation)

\textbf{Step 2:} Substitute into the utility constraint:
\[
u^* = \min\{4x_1, x_2\} = \min\{4x_1, 4x_1\} = 4x_1
\]

\textbf{Step 3:} Solve for $x_1$:
\[
x_1 = \frac{u^*}{4}
\]

\textbf{Step 4:} Use the kink condition to find $x_2$:
\[
x_2 = 4x_1 = 4 \cdot \frac{u^*}{4} = u^*
\]

\subsection{Final Answer}

\[
\boxed{h_1(p_1, p_2, u^*) = \frac{u^*}{4}}
\]
\[
\boxed{h_2(p_1, p_2, u^*) = u^*}
\]

This is answer \textbf{(b)}: $\left(\frac{u^*}{4}, u^*\right)$

\subsection{Properties}

\begin{itemize}
    \item Depends on \textbf{target utility $u^*$} (not income!)
    \item For perfect complements, \textbf{independent of prices} (the ratio is always fixed)
    \item Also called ``compensated demand''
    \item Theoretical construct, not directly observable
\end{itemize}

\section{Key Difference: Prices Don't Affect Hicksian for Perfect Complements}

Notice something surprising: Hicksian demand $\left(\frac{u^*}{4}, u^*\right)$ doesn't depend on $p_1$ or $p_2$!

\textbf{Why?} With perfect complements, you \textbf{must} consume in the fixed ratio $x_2 = 4x_1$ no matter what. If you want utility $u^*$, you need exactly:
\begin{itemize}
    \item $\frac{u^*}{4}$ units of good 1
    \item $u^*$ units of good 2
\end{itemize}

There's no substitution possible. Prices only affect the \textbf{total cost}, not the \textbf{quantities}.

\section{Comparison Table}

\begin{center}
\begin{tabular}{|l|c|c|}
\hline
\textbf{Feature} & \textbf{Marshallian} & \textbf{Hicksian} \\
\hline
Problem & Max $u$ s.t. $p \cdot x = m$ & Min $p \cdot x$ s.t. $u = u^*$ \\
\hline
Fixed variable & Income $m$ & Utility $u^*$ \\
\hline
Kink condition & $x_2 = 4x_1$ & $x_2 = 4x_1$ \\
\hline
$x_1$ demand & $\dfrac{m}{p_1 + 4p_2}$ & $\dfrac{u^*}{4}$ \\
\hline
$x_2$ demand & $\dfrac{4m}{p_1 + 4p_2}$ & $u^*$ \\
\hline
Depends on prices? & Yes & No (for perfect compl.) \\
\hline
Depends on income? & Yes & No \\
\hline
Depends on utility? & No (achieved, not fixed) & Yes (fixed) \\
\hline
\end{tabular}
\end{center}

\section{Numerical Example}

Let $p_1 = 1$, $p_2 = 2$, and $m = 18$.

\subsection{Marshallian Approach}

Find the optimal bundle given income $m = 18$:
\begin{align*}
x_1^* &= \frac{18}{1 + 4(2)} = \frac{18}{9} = 2 \\
x_2^* &= \frac{4(18)}{1 + 4(2)} = \frac{72}{9} = 8
\end{align*}

Check: $p_1 x_1 + p_2 x_2 = 1(2) + 2(8) = 18$ \checkmark

Achieved utility: $u^* = \min\{4(2), 8\} = \min\{8, 8\} = 8$

\subsection{Hicksian Approach}

Now suppose we want to achieve utility $u^* = 8$ as cheaply as possible:
\begin{align*}
h_1 &= \frac{8}{4} = 2 \\
h_2 &= 8
\end{align*}

Total expenditure: $e(p_1, p_2, u^*) = 1(2) + 2(8) = 18$

\subsection{Duality}

Notice: Both approaches give the \textbf{same bundle} $(2, 8)$!

This demonstrates the duality relationship:
\[
x^*(p, m) = h(p, u^*) \quad \text{where } u^* = u(x^*(p, m))
\]

The Marshallian demand at income $m$ equals the Hicksian demand at the utility level achieved with income $m$.

\section{Why Answer (b) is Correct}

The question asks for Hicksian demand with utility $u^*$. Using the constraint $\min\{4x_1, x_2\} = u^*$ and the kink condition $4x_1 = x_2$:

\begin{align*}
u^* &= 4x_1 = x_2 \\
\Rightarrow x_1 &= \frac{u^*}{4} \\
\Rightarrow x_2 &= u^*
\end{align*}

Therefore, the Hicksian demand bundle is: $\boxed{\left(\frac{u^*}{4}, u^*\right)}$ which is answer \textbf{(b)}.

\section{Common Mistakes}

\begin{enumerate}
    \item \textbf{Confusing Marshallian with Hicksian:} Marshallian depends on $m$, Hicksian depends on $u^*$

    \item \textbf{Getting the ratio backwards:} The constraint is $u = \min\{4x_1, x_2\}$, so we need $4x_1 = x_2$, which means $x_1 = \frac{x_2}{4}$ or equivalently $x_1 = \frac{u^*}{4}$

    \item \textbf{Forgetting prices don't matter:} For perfect complements, Hicksian demand doesn't depend on prices because the ratio is fixed
\end{enumerate}

\section{Summary}

\begin{itemize}
    \item \textbf{Perfect complements} require fixed consumption ratio: $x_2 = 4x_1$

    \item \textbf{Marshallian demand}: Given income $m$, buy $\left(\frac{m}{p_1 + 4p_2}, \frac{4m}{p_1 + 4p_2}\right)$

    \item \textbf{Hicksian demand}: To achieve utility $u^*$, buy $\left(\frac{u^*}{4}, u^*\right)$

    \item For perfect complements, Hicksian demand is \textbf{independent of prices}

    \item The two demands are related through duality: same bundle at corresponding income/utility levels
\end{itemize}

\end{document}
